% Part: applied-modal-logics
% Chapter: temporal-logic
% Section: possible-histories

\documentclass[../../../include/open-logic-section]{subfiles}

\begin{document}

\olfileid{aml}{tl}{poss}

\olsection{Riwayat Kang Mungkin}

Model relasional logika temporal kang wis kita gunakake iku luwes
banget, amarga kita ora kudu netepake watesan tumrap relasi
aksesibilitas. Mula model temporal bisa nyabang ing kala kepungkur
lan mangsa ngarep. Nanging, kita bisa nimbang pangertèn percabangan
kang luwih ``modal'', yaiku ndeleng urutan kedadeyan minangka
riwayat kang mungkin. Iki ora mesthi mbutuhake owah-owahan basa,
senajan kita uga bisa nambahake operator modal ``lumrah'' $\Box$
lan~$\Diamond$, sarta relasi aksesibilitas epistemik kanggo makili
owah-owahan kawruh para agen saka wektu menyang wektu.

\begin{defn}
  \emph{Model riwayat kang mungkin} kanggo basa temporal yaiku tripel
  $\mModel{M} = \tuple{T, C, V}$, kanthi
  \begin{enumerate}
    \item $T$ iku himpunan ora kosong, kang diinterpretasikake
      minangka kahanan ing wektu.
    \item $C$ iku himpunan lintasan komputasi, utawa \emph{riwayat
      kang mungkin} sawijining sistem. Tegese, $C$~iku himpunan
      urutan~$\sigma$ saka kahanan $s_1$, $s_2$,~$s_3$, \dots, kang
      saben $s_i \in T$.
    \item $V$ iku fungsi kang netepake marang saben !!{propositional
      variable}~$p$ sawijining himpunan~$V(p)$ saka titik wektu.
  \end{enumerate}
  Supaya luwih prasaja, kita uga lumrahe nganggep yen sawijining
  riwayat ana ing~$C$, kabeh sufikse uga ana ing kono. Contone, yen
  $s_1$, $s_2$,~$s_3$ minangka urutan ing~$C$, mula $s_2$, $s_3$
  lan~$s_3$ uga ana ing kono. Yen rong kahanan $s_i$ lan~$s_j$ katon
  ing urutan~$\sigma$, kita kandha $s_i \prec_\sigma s_j$ yen $i < j$.
  Yen $t \in V(p)$, kita kandha $p$~\emph{bener ing}~$t$.
\end{defn}

Owah-owahan kang wigati yaiku nalika kita ngevaluasi bebener
!!a{formula} ing titik wektu~$t$ sajroning model~$\mModel M$,
kita nindakake iku relatif marang riwayat~$\sigma$, kang ngemot $t$
minangka kahanan. Nanging, semantik kanggo !!{propositional
variable}s lan konektif kang fungsional miturut bebener ora perlu
diowahi. Kabeh iku tetep kaya ing \olref[sem]{defn:tmodels}, amarga
ora ana kang nyebut~$\sigma$. Saiki operator mangsa ngarep~$\Ftemp$
kita tegesi maneh lan operator~$\Diamond$ kita tambahake relatif
marang riwayat kasebut.

\begin{defn}\ollabel{defn:phmodels}
  \emph{Bebener !!a{formula}~$!A$ ing~$t, \sigma$} sajroning~$\mModel M$, kanthi lambang:
  $\mSat{M}{!A}[t, \sigma]$:
  \begin{enumerate}
  \item\ollabel{defn:sub:phmodels-f} \indcase{!A}{\Ftemp
    !B}{$\mSat{M}{\indfrm}[t, \sigma]$ yen lan mung yen $\mSat{M}{!B}[t', \sigma]$
    kanggo sawijining $t' \in T$ kang nyukupi $t \prec_\sigma t'$.}
  \item\ollabel{defn:sub:phmodels-diamond} \indcase{!A}{\Diamond
    !B}{$\mSat{M}{\indfrm}[t, \sigma]$ yen lan mung yen $\mSat{M}{!B}[t, \sigma']$
    kanggo sawijining $\sigma' \in C$ kang ngemot $t$.}
  \end{enumerate} 
\end{defn}

Operator temporal lan modal liyane bisa ditegesi kanthi cara kang
padha. Saiki kita bisa makili pratelan kang nggabungake kala lan
modalitas. Contone, ``$p$~ora bakal kelakon, nanging bisa wae
kelakon'' bisa dilambangake nganggo rumus
$\lnot \Ftemp p \land \Diamond \Ftemp p$. Rumus iki bener ing
titik lan riwayat kang $p$ ora dadi bener ing kahanan peneruse,
nanging ana riwayat alternatif kang $p$ bakal dadi bener.

\end{document}
